% ============================================================
% f2014_beijing_arts_q04_flowchart.tex
% 2014 北京卷（文科）第 4 题：程序框图
% 来源 PDF：普通高考/2014/2014北京文.pdf
% 原页：第 1 页（A3），印面 863，左栏第 4 题题干下方
% 结构：开始 → k=0, S=0 → 判断 k<3；是则 S=S+2^k、k=k+1 并返回判断；
%       否则向右输出 S，再向下结束。
% ============================================================
\begin{tikzpicture}[>=Stealth, line width=0.7pt, font=\small]
  % 开始
  \draw[rounded corners=6pt] (-0.95,0.45) rectangle (0.95,-0.45);
  \node at (0,0) {开始};
  \draw[->] (0,-0.45) -- (0,-1.05);

  % 初始化
  \draw (-1.75,-1.05) rectangle (1.75,-2.05);
  \node at (0,-1.55) {$k=0$，$S=0$};
  \draw (0,-2.05) -- (0,-2.35);
  \draw[->] (0,-2.35) -- (0,-2.65);

  % 判断菱形
  \draw (0,-2.65) -- (1.55,-3.45) -- (0,-4.25) -- (-1.55,-3.45) -- cycle;
  \node at (0,-3.45) {$k<3$};

  % 是：循环体
  \draw[->] (0,-4.25) -- (0,-4.95);
  \node at (-0.22,-4.62) {是};
  \draw (-1.75,-4.95) rectangle (1.75,-5.95);
  \node at (0,-5.45) {$S=S+2^k$};
  \draw[->] (0,-5.95) -- (0,-6.35);
  \draw (-1.75,-6.35) rectangle (1.75,-7.35);
  \node at (0,-6.85) {$k=k+1$};

  % 回边：由 k=k+1 左侧绕回判断框之前的竖线
  \draw (-1.75,-6.85) -- (-2.85,-6.85) -- (-2.85,-2.35);
  \draw[->] (-2.85,-2.35) -- (0,-2.35);

  % 否：输出
  \draw[->] (1.55,-3.45) -- (3.35,-3.45);
  \node at (2.35,-3.20) {否};
  \draw[->] (3.35,-3.45) -- (3.35,-3.98);
  \draw (1.95,-4.68) -- (4.45,-4.68) -- (4.15,-3.98) -- (1.65,-3.98) -- cycle;
  \node at (3.05,-4.33) {输出 $S$};
  \draw[->] (3.05,-4.68) -- (3.05,-5.12);

  % 结束
  \draw[rounded corners=6pt] (2.15,-5.12) rectangle (3.95,-6.02);
  \node at (3.05,-5.57) {结束};
\end{tikzpicture}
